% Ad Atticum 4.4A --- headnote
% ad-atticum-04-04a, April or May 56 BC, written at Antium

\textit{Cicero to Atticus, written at Antium in April or May 56
BC. The famous library letter, and the most intimate glimpse
in all the surviving correspondence of how a Roman gentleman's
library was actually put back on its feet. Cicero is at his
seaside villa at Antium with the books that survived the
Clodian sack of the previous year, and Tyrannio --- the Greek
grammarian Tyrannion of Amisus, brought to Rome by Lucullus
after the Mithridatic war and now the leading scholar in the
city --- has been arranging the rolls. The reliquiae, Cicero
discovers, are much better than he had thought.}

\textit{The technical request that follows is what makes the
letter unique. He asks Atticus to send two of his trained
book-slaves: glutinatores, who paste the papyrus sheets into
rolls and mend torn ones, and ad cetera administris, hands for
the rest of the work. They are to bring strips of vellum
(membranula) on which to write the title-labels --- the
\textit{[Greek: sittybai]} that hung from the projecting end of a
roll so a reader could find it on the shelf without unrolling
each. The light teasing of Atticus's Greek is part of the
intimacy: ``what you Greeks, I believe, call sittybai.'' Behind
this paragraph stands the whole machinery of the late-Republican
private library.}

\textit{The letter closes with two sharper notes from Roman
public life: Atticus has bought a troop of gladiators, and
Cicero (who hears they fight wonderfully) teases him that, had
he hired them out for the two shows now in prospect, the
purchase would have paid for itself. And the steady refrain of
the letters of these months: come, and bring Pilia. Tullia
wants it.}
